\subsection*{Бинарный поиск}
\Task Дан отсортированный массив A[1 : n]. Надо проверить, есть ли в нем элемент X. 

\Alg
Обозначим за $l = 1$ и $r = n$ левую и правую границы
рассматриваемого подмассива.
\begin{enumerate}
    \item Если диапазон валиден, то проверим, $ A[ \left\lfloor \frac{l+r}{2} \right\rfloor] \geq X$. 
    Иначе вернем границу.
    \item Если полученный результат истинен, то запускаем рекурсивно
поиск в левой половине массива.
    \item Иначе — в правой половине массива.
\end{enumerate}

\textit{Асимптотика.}
\begin{enumerate}
    \item Сравнивает средний элемент с X за O(1).
    \item Решает, продолжать в левой или правой половине за O(1).
    \item Запускается рекурсивно, уменьшая область поиска в два раза.
\end{enumerate}
Результат будет получен, когда длина диапазона (далее n) станет
равной 1. 
А значит всего шагов $\log_2 n$
$$O(log_2\,n) = O(log\,n)$$

$T(N) = T(n/2) + O(1)$. По Мастер-теореме $T(N) = O(log\,n)$

\subsection*{Бинарный поиск по ответу}
Пусть имеется монотонный предикат \( P(n) \), то есть
$$
\exists N_0 : 
\begin{cases}
P(n) = 0, & n < N_0 \\
P(n) = 1, & n \geq N_0
\end{cases}
$$

И перед нами стоит задача отыскать это \( N_0 \). Например, задача проверки наличия элемента \( X \) в отсортированном массиве. В данном случае предикат будет звучать как 

$$P(i) = 1 \iff a[i] \geq X.$$


\Alg
Пусть \( N_0 \) заведомо лежит в каком-то диапазоне \([l, r]\). Тогда алгоритм следующий:
\begin{enumerate}
    \item Если диапазон валиден, то вычислим \( P\left( \left\lfloor \frac{l+r}{2} \right\rfloor \right) \). Иначе вернем границу.
    \item Если полученный результат истинен, то \( N_0 \in [l, \left\lfloor \frac{l+r}{2} \right\rfloor] \). Запускаем рекурсивно поиск в левой половине массива.
    \item Иначе \( N_0 \in \left[\left\lfloor \frac{l+r}{2} \right\rfloor, r\right] \). Запускаем рекурсивно поиск в правой половине массива.
\end{enumerate}
\textit{Асимптотика.} \( O(\log n) \).
\pagebreak